The deeper function of vision is not merely to output labels, but to recover a stable and actionable model of the external world from incomplete and noisy visual input. Inspired by psychophysics and systems neuroscience, we treat a vision-language model (VLM) as a black-box observer and probe its \emph{externalized scene belief} through dense ordinal spatial questions. This is motivated by a simple observation from ordinal embedding: if sufficiently many ordinal relations are correct, a scene can be approximately reconstructed up to a similarity transform; asking a VLM enough QRR and TRR questions is therefore a way of testing what kind of scene the model effectively ``has in mind.'' Our current experiments on 540 completed scenes and over 200K QA pairs show that scene-level spatial understanding remains weak: single-view QRR reaches only about 70\% accuracy, multi-view QRR drops toward chance, TRR stays near chance, and neither supervised fine-tuning nor reinforcement learning on these relational QA pairs yields substantial gains. Ordinary-Bench thus shifts evaluation from ``how many spatial questions did the model answer correctly?'' to ``does the model support a coherent, reconstructable, and alignable spatial world model at all?'' and opens a path toward vectorized scene reconstruction, localization, and structured scene memory from ordinal relations.
